%==============================================================================
% CITS3402/CITS5507 Assignment 1 — Report template
%
% Build:  pdflatex report.tex && pdflatex report.tex
%         (twice, so the ToC/refs settle)
%
% Word count (target 1000 +/- 100, body text only):
%         texcount -inc -sum -1 report.tex
% The \wordcountnote command below prints a reminder; delete it before you
% submit.
%==============================================================================

\documentclass[11pt,a4paper]{article}

%------------------------------------------------------------------ packages --
\usepackage[T1]{fontenc}
\usepackage[utf8]{inputenc}
\usepackage{lmodern}
\usepackage[margin=2.5cm]{geometry}
\usepackage{microtype}

\usepackage{amsmath,amssymb}
\usepackage{booktabs}          % proper table rules
\usepackage{siunitx}           % \SI{1.23}{\second}, aligned numeric columns
\usepackage{graphicx}
\usepackage{caption}
\usepackage{subcaption}
\usepackage{listings}
\usepackage{xcolor}
\usepackage{pgfplots}
\pgfplotsset{compat=1.18}
\usepackage[hidelinks]{hyperref}
\usepackage{cleveref}

%--------------------------------------------------------------- code listing --
\definecolor{codegray}{gray}{0.95}
\lstdefinestyle{c}{
  language=C,
  basicstyle=\ttfamily\footnotesize,
  keywordstyle=\color{blue!60!black},
  commentstyle=\color{green!40!black}\itshape,
  stringstyle=\color{red!60!black},
  backgroundcolor=\color{codegray},
  numbers=left, numberstyle=\tiny\color{gray}, numbersep=6pt,
  breaklines=true, showstringspaces=false, frame=none,
  morekeywords={\#pragma,omp,parallel,for,critical,atomic,uint64_t,uint32_t,size_t}
}
\lstset{style=c}

%---------------------------------------------------------------- convenience --
\newcommand{\toyhash}{\texttt{toy\_hash}}
\newcommand{\wordcountnote}[1]{\marginpar{\scriptsize\textcolor{gray}{#1}}}

\sisetup{round-mode=places, round-precision=2, table-format=3.2}

%--------------------------------------------------------------------- title --
\Huge \title{\vspace{-2em}CITS3402 - High Performance Computing\\[0.3em]
\LARGE Assignment 1: Report  \\[1.0em]
\large The University of Western Australia}

\date{Semester 2, 2026}

% Group of two: list both. Solo: delete the second block.
\author{
  James Oakey \quad (\texttt{22709404})\\
  \and
  Brent Tankink\quad (\texttt{23842541})
}


\begin{document}
\maketitle

%==============================================================================
\section{Introduction}
%==============================================================================

This document outlines how the puzzling results of the Birthday Paradox can be
exploited in order to brute force hash collisions. The goal is to alter the 
'nonce' values in two PDF files, each with differing contents, in such a way 
that the resultant hash is identical. 

This collision is produced through the combination of nonce variation and the 
Birthday Attack principle, which proves that collision probability increases 
rapidly with less trials than is intuitively expected. 

%==============================================================================
\section{Birthday-Attack Algorithm}
\label{sec:bday}
%==============================================================================

The implemented search strategy uses a Hash-Table data structure to store and 
search through the calculated hashes. Rather than comparing every possible 
nonce produced from one PDF file to the other, each calculated (hash+file) 
combination is checked against those already in the table. Essentially, rather 
than checking whether a match occurs with one specific hash, the program checks 
whether a match occurs between any two hashes. 

The $\texttt{toy\_hash()}$ algorithm produces 48-bit hashes, meaning there are $2^{48}$ 
possible hashes. Using the Birthday Paradox, the expected number of trials for 
a 50\% collision probability follows $2^{\sqrt{48}}=2^{24}$, meaning 
only ~16-17 million calculations are required to achieve this 50\% probability.
To find an upper bound for the calculation process, the following formula can 
be used: $P(\text{No Collision})\approx e^{-\frac{n^2}{2H}}$, where $n$ is the 
number of trials required to achieve $P$ probability of no collision when there 
are $H$ possible outputs. For a 99\% chance of producing a collision:

$$0.01=e^{-\frac{n^2}{2\times2^{48}}}\quad\rightarrow\quad n\approx51,000,000$$

For the stated reason above, the maximum number of iterations is capped at 
50,000,000, as this will give a very high probability of collision and won't 
require a huge initial memory allocation.

To store generated hashes, a hash table has been implemented using linear 
probing. Each entry, $\texttt{HashEntry}$, contains a 64-bit tag and a 64-bit nonce, 
making each entry 16 bytes. The $\texttt{tag}$ attribute combines the PDF file number 
(1 or 2) and the 48-bit file hash, making efficient use of space and allowing 
the program to distinguish hashes from PDF1 and PDF2. This table is initialised 
using $\texttt{calloc}$, making each entry 0. Whenever a new hash is generated, the 
lowest 27 bits of the hash value are used to calculate the insertion index in 
the hash table. If slot is already allocated, linear probing is performed to 
find the next free slot. The produced hashes and file identifiers are then 
checked against previous ones, allowing the collision detection between the two 
PDFs. 

The hash table created contains $2^{27}$ 16-byte entries, meaning the table 
uses just over 2GB. 

%==============================================================================
\section{Parallelisation with OpenMP}
%==============================================================================

The serial implementation works by assigning a different nonce to each PDF in 
every loop iteration, then calculating the hash for both files. These hashes 
are then checked against the hash table, if either of them match a hash which 
was produced by the other PDF, then a collision is found and the collision pair 
is returned (if no collision is found then the two are inserted in the table). 
The table insertion is done through a function named $\texttt{insert\_or\_match()}$, which 
performs the lookup and probing outlined in the birthday attack section. The 
basic process is as follows:

\begin{center}
    \begin{tabular}{l}
        $\mid$ 1. Set and insert nonces into PDF 1 and 2\\
        $\mid$ 2. Calculate hash of each PDF\\
        $\mid$ 3. Check hash table\\
        $\hookrightarrow$ 3a. If match: save nonces and stop\\
        $\hookrightarrow$ 3b. If not: insert hashes into table and continue
    \end{tabular}
\end{center}

The parallel implementation works by splitting the nonce generation and the PDF 
hashing across multiple threads in a $\texttt{parallel}$ region. The main $\texttt{for}$ loop is split into a parallel 
region with $\texttt{schedule(static)}$, as each iteration performs a 
similar amount of work. 

Each thread creates its own copies of the two PDFs so that the $\texttt{set\_nonce()}$ 
function doesn't continuously overwrite each thread's contents (which could 
cause incorrect hash calculations), eliminating a race condition. The hash 
table is shared between all threads, so the lookup and insertion operations are 
isolated inside a critical region. This way, there will be no data races 
occurring in the shared table. Without this critical region, two threads could 
simultaneously try to insert a hash value in the same hash table location, 
making the table state inconsistent. The critical region ensures that only 1 
thread accesses the table at a time. 

The computationally expensive operations are placed outside of this critical 
region, with those being $\texttt{set\_nonce()}$ and $\texttt{toy\_hash()}$. This minimises the 
size of the critical region and the time that threads are prevented from 
running concurrently. The $\texttt{found}$ variable is also shared, which in the serial 
implementation checks whether a collision has occurred. In the parallel code, 
this variable needs to be synchronised so that all threads know when a 
collision has occurred. Each iteration $\texttt{atomic read}$'s the $\texttt{found}$ variable before 
entering the table checker, which makes sure that any collision is detected. 
Upon entering the critical region, this $\texttt{found}$ value is checked again, as 
there is a possibility that a collision has occurred within the small time 
frame between the initial check and the critical region.

There are a couple of trade-offs with this approach, and these are discussed 
further in the following section.

%==============================================================================
\section{Memory Usage \& Trade-Offs}
%==============================================================================

As stated in the \hyperref[sec:bday]{Birthday-attack algorithm}, the hash table 
size is set as $2^{27}\times16 \text{ Bytes}\approx2\text{GB}$ bytes, capping 
the memory at just over 2 gigabytes. This was picked to ensure that the maximum 
50,000,000 trials (for each PDF) can be stored while having some space 
remaining to minimise hash table probing collisions. While dynamic hash table 
size increasing would have been more memory efficient, this would have created 
a much longer running time, hence it was decided to allocate the large initial 
hash table.

A particular trade-off made in the use of a critical region instead of a multi-
partitioned table as the critical region approach has it's speed limited by the 
shared table. Increasing the thread count will increase the concurrent nonce 
and hash calculations, but the critical region means only one thread can access 
the table at a time, bottlenecking the program. An alternative to this could be 
a partitioned hash table (mentioned in the task brief), where different threads 
can access different partitions side by side. The partitioned approach would 
have produced a speed-up, however it would likely have been much more difficult 
to implement. The possibility for increased speed was traded for simplicity.

%==============================================================================
\section{Performance Evaluation}
%==============================================================================

\subsection{Methodology}




\subsection{Speedup and scalability}

\subsection{Discussion}


%==============================================================================
\section{Conclusion}
%==============================================================================

%==============================================================================
\appendix
\section{Reproducing the results}

\end{document}
